\subsection{附加于有限生成模的除子类}

我们保持第 2 至第 6 节的记号与假设。回忆：C(A)（或简作 C）表示 A 的除子类群，即 D(A) 关于主除子所成子群的商。对每个除子 \(d \in D\) ，我们用 \(c(d)\) 表示它在 C 中的类。

命题 15. 设 M 是有限生成 A-模。存在 M 的自由子模 L，使 M/L 是扭模，且 C 的元素 \(c(\chi(M/L))\) 与所选的自由子模 L 无关。

我们写 \(S = A - \{0\}\) 并令 \(V = S^{-1}M = M \otimes_{A} K\) ；若 \(n\) 是 \(V\) 在 \(K\) 上的秩，则存在 \(n\) 个元素 \(e_i (1 \leqslant i \leqslant n)\) 属于 \(M\) ，它们的典范像在 \(V\) 中构成 \(V\) 的基；这些元素在 \(M\) 中显然线性无关，从而生成自由子模 \(L\) ，它含于 \(M\) ，满足 \(S^{-1}(M/L) = S^{-1}M/S^{-1}L = 0\) ，因此 \(M/L\) 是扭模。

现在设 \(L_{1}\) 是 M 的另一个秩为 n 的自由子模。~由于 \(S^{-1}L = S^{-1}L_{1}\) ，存在 \(s \in S\) 使 \(sL_{1} \subset L\) ；因此我们只需限于证明：若 \(L_{1} \subset L_{2}\) 是 M 的两个秩为 n 的自由子模，则

\[
c (\chi (\mathbf {M} / \mathbf {L} _ {1})) = c (\chi (\mathbf {M} / \mathbf {L} _ {2})).
\]

现在，\(\chi (\mathbf{M} / \mathbf{L}_1) = \chi (\mathbf{M} / \mathbf{L}_2) + \chi (\mathbf{L}_2 / \mathbf{L}_1)\) ，且由第 6 节命题 14 的推论可知 \(\chi (\mathbf{L}_2 / \mathbf{L}_1)\) 是主除子，因此

\[
c (\chi (\mathbf {L} _ {2} / \mathbf {L} _ {1})) = 0.
\]

元素 \(c(\chi(\mathbf{M} / \mathbf{L}))\) 以下记作 \(-c(\mathbf{M})\) ；我们称 \(c(\mathbf{M})\) 为附加于 \(\mathbf{M}\) 的除子类。

命题 16. (i) 设 \(0 \to \mathbf{M}_1 \xrightarrow{f} \mathbf{M}_2 \xrightarrow{g} \mathbf{M}_3 \to 0\) 是有限生成 A-模的正合列。则

\[
c (\mathbf {M _ {2}}) = c (\mathbf {M _ {1}}) + c (\mathbf {M _ {3}}).
\]

\begin{enumerate}
\def\labelenumi{(\roman{enumi})}
\setcounter{enumi}{1}
\item
  若存在从 \(\mathbf{M}_1\) 到 \(\mathbf{M}_2\) 的伪同构，则 \(c(\mathbf{M}_1) = c(\mathbf{M}_2)\) 。
\item
  若 \(\mathbf{T}\) 是扭模，则 \(c(\mathbf{T}) = -c(\chi(\mathbf{T}))\) 。
\item
  若 \(\mathfrak{a} \neq 0\) 是 \(\mathbf{K}\) 的分式理想，则
\end{enumerate}

\[
c (\mathfrak {a}) = c (\operatorname{div} (\mathfrak {a})).
\]

\begin{enumerate}
\def\labelenumi{(\alph{enumi})}
\setcounter{enumi}{21}
\tightlist
\item
  若 \(\mathbf{L}\) 是自由 A-模，则 \(c(\mathbf{L}) = 0\) 。
\end{enumerate}

为证 (i)，考虑自由子模 \(\mathbf{L}_1\) （相应地，\(\mathbf{L}_3\) ），它含于 \(\mathbf{M}_1\) （相应地，\(\mathbf{M}_3\) ），使 \(\mathbf{M}_1 / \mathbf{L}_1\) （相应地，\(\mathbf{M}_3 / \mathbf{L}_3\) ）是扭模。由于 \(\mathbf{L}_3\) 自由且 \(g\) 是满射，在 \(\overline{g}^{-1} (\mathbf{L}_3)\) 中存在一个自由补 \(\mathbf{L}_{23}\) （ \(\operatorname {Ker}(g)\) 的），它与 \(\mathbf{L}_3\) 同构 （《代数》第 II 章，§ 1，第 11 节，命题 21）；但 \(\operatorname {Ker}(g) = f(\mathbf{M}_1)\)

它包含 \(f(\mathbf{L}_1) = \mathbf{L}_{12}\) ，后者因 \(f\) 是单射而是自由的。\(\mathbf{L}_2 = \mathbf{L}_{12} + \mathbf{L}_{23}\) 这个和是直和，因此 \(\mathbf{L}_2\) 是 \(\mathbf{M}_2\) 的自由子模。此外还有交换图：

\[
\begin{array}{c} 0 \longrightarrow \mathrm{L} _ {1} \longrightarrow \mathrm{L} _ {2} \longrightarrow \mathrm{L} _ {3} \longrightarrow 0 \\ \Biggl \downarrow \qquad \qquad \qquad \Biggl \downarrow \qquad \qquad \Biggl \downarrow \\ 0 \longrightarrow \mathrm{M} _ {1} \xrightarrow [ f ]{} \mathrm{M} _ {2} \xrightarrow [ g ]{} \mathrm{M} _ {3} \longrightarrow 0 \end{array}
\]

其中各行正合，各竖直箭均为单射。因此我们由蛇形图（第 I 章，§ 1，第 4 节，命题 2）得到正合列：

\[
0 \rightarrow \mathrm{M} _ {1} / \mathrm{L} _ {1} \rightarrow \mathrm{M} _ {2} / \mathrm{L} _ {2} \rightarrow \mathrm{M} _ {3} / \mathrm{L} _ {3} \rightarrow 0.
\]

由于 \(\mathbf{M}_1 / \mathbf{L}_1\) 与 \(\mathbf{M}_3 / \mathbf{L}_3\) 都是扭模，这一正合列首先表明 \(\mathbf{M}_2 / \mathbf{L}_2\) 亦然，继而由第 5 节命题 10 得：

\[
\chi (\mathrm{M} _ {2} / \mathrm{L} _ {2}) = \chi (\mathrm{M} _ {1} / \mathrm{L} _ {1}) + \chi (\mathrm{M} _ {3} / \mathrm{L} _ {3})
\]

这就证明了 (i)。

断言 (iii) 与 (v) 由定义立得。我们证明 (ii)。于是设 \(f \colon \mathbf{M}_1 \to \mathbf{M}_2\) 是伪同构，\(\mathbf{L}_1\) 是 \(\mathbf{M}_1\) 的自由子模，使 \(\mathbf{M}_1 / \mathbf{L}_1\) 是扭模。我们置 \(\mathbf{L}_2 = f(\mathbf{L}_1)\) ；由于 \(\operatorname{Ker}(f)\) 是伪零的，它是扭模，故 \(\operatorname{Ker}(f) \cap \mathbf{L}_1 = 0\) ，从而 \(\mathbf{L}_2\) 是自由的。设 \(\bar{f} \colon \mathbf{M}_1 / \mathbf{L}_1 \to \mathbf{M}_2 / \mathbf{L}_2\) 是由 \(f\) 取商导出的同态；\(\operatorname{Ker}(\bar{f})\) 同构于 \(\operatorname{Ker}(f)\) ，\(\operatorname{Coker}(\bar{f})\) 同构于 \(\operatorname{Coker}(f)\) ，故 \(\bar{f}\) 是伪同构；此外

\[
\operatorname{Coker} (\bar {f}) = \mathrm{M} _ {2} / f (\mathrm{M} _ {1})
\]

是扭模，\(f(\mathbf{M}_1) / \mathbf{L}_2 = \bar{f} (\mathbf{M}_1 / \mathbf{L}_1)\) 亦然，故 \(\mathbf{M}_2 / \mathbf{L}_2\) 是扭模，于是由第 5 节命题 10 (ii) 得

\[
\chi (\mathrm{M} _ {1} / \mathrm{L} _ {1}) = \chi (\mathrm{M} _ {2} / \mathrm{L} _ {2}).
\]

最后尚须证明 (iv)。设 \(x \in \mathbf{K}^*\) 满足 \(a \subset xA\) 。考虑正合列 \(0 \to a \to xA \to xA / a \to 0\) ，我们得到

\[
c (\mathfrak {a}) = c (x \mathrm{A}) - c (x \mathrm{A} / \mathfrak {a}) = - c (x \mathrm{A} / \mathfrak {a})
\]

这是依据 (i) 与 (v)。但 \(x\mathrm{A} / \mathfrak{a}\) 同构于 \(\mathrm{A} / x^{-1}\mathfrak{a}\) ，于是依据 (iii)，

\[
c (x \mathrm{A} / \mathrm{a}) = - c (\chi (\mathrm{A} / x ^ {- 1} \mathrm{a})) = - c (\operatorname{div} (x ^ {- 1} \mathrm{a})) = - c (\operatorname{div} (\mathrm{a}))
\]

（第 5 节，命题 12）。证毕。

当 M 是 V 关于 A 的格时，\(\chi(\mathrm{M}/\mathrm{L}) = -\chi(\mathrm{M}, \mathrm{L})\) （第 6 节，命题 14）；设 \((e_{i})_{1 \leqslant i \leqslant n}\) 是 L 的基，\(e = e_{1} \wedge e_{2} \wedge \cdots \wedge e_{n}\) ，且

\(M_{w} = a.e\) （第 6 节的记号）；则 \(\chi(M, L) = \text{div}(a)\) ，从而 \(c(M) = c(\text{div}(a))\) ，这推广了命题 16 (v)。

推论 1. 设 \(0 \to \mathbf{M}_n \xrightarrow{u} \mathbf{M}_{n-1} \to \cdots \to \mathbf{M}_0 \to 0\) 是有限生成 A-模的正合列。则

\[
\sum_ {i = 0} ^ {n} (- 1) ^ {i} c (\mathbf {M} _ {i}) = 0.
\]

我们对 \(n\) 作归纳，\(n = 2\) 的情形即命题 16 (i)。若 \(\mathbf{M}_{n-1}' = \operatorname{Coker}(u)\) ，则有两个正合列：

\[
0 \rightarrow \mathrm{M} _ {n} \rightarrow \mathrm{M} _ {n - 1} \rightarrow \mathrm{M} _ {n - 1} ^ {\prime} \rightarrow 0
\]

\[
0 \rightarrow \mathrm{M} _ {n - 1} ^ {\prime} \rightarrow \mathrm{M} _ {n - 2} \rightarrow \dots \rightarrow \mathrm{M} _ {0} \rightarrow 0.
\]

第一个正合列表明 \(\mathbf{M}_{n - 1}^{\prime}\) 是有限生成的，且归纳假设给出

\[
(- 1) ^ {n - 1} c (\mathbf {M} _ {n - 1} ^ {\prime}) + \sum_ {i = 0} ^ {n - 2} (- 1) ^ {i} c (\mathbf {M} _ {i}) = 0
\]

且

\[
c (\mathbf {M} _ {n - 1} ^ {\prime}) = c (\mathbf {M} _ {n - 1}) - c (\mathbf {M} _ {n}),
\]

由此得推论。

正合列

\[
0 \rightarrow \mathrm{L} _ {n} \rightarrow \mathrm{L} _ {n - 1} \rightarrow \dots \rightarrow \mathrm{L} _ {0} \rightarrow \mathrm{E} \rightarrow 0
\]

其中诸 \(L_{i}\) （ \(0 \leqslant i \leqslant n\) ）是有限生成的自由 A-模，这样的正合列称为 A-模 E 的有限自由分解。

推论 2. 若 A 的除性分式理想 \(\mathfrak{a} \neq 0\) 容许一个有限自由分解，则它是主理想。

事实上，我们对 \(\mathfrak{a}\) 的一个有限自由分解应用推论 1：

\[
0 \rightarrow \mathrm{L} _ {n} \rightarrow \mathrm{L} _ {n - 1} \rightarrow \dots \rightarrow \mathrm{L} _ {0} \rightarrow \mathfrak {a} \rightarrow 0.
\]

由命题 16 (v)，\(c(a) = 0\) ，从而由命题 16 (iv)，\(\text{div}(a)\) 是主理想；由于假设 \(a\) 是除性的，故它是主理想（§ 1，第 1 节）。

推论 3. 若 A 的每个 \(\neq 0\) 除性理想都容许有限自由分解，则 A 是因子分解整环。

这是推论 2 与 § 3 第 1 节定义 1 的直接推论。

* 我们以后将看到，正则局部环满足推论

3 的假设，从而是因子分解整环。*

若 M 是有限生成 A-模，我们将用 \(r(\mathbf{M})\) 表示它的秩（回忆它是 \(\mathbf{M}_{(\mathbf{K})} = \mathbf{M} \otimes_{\mathbf{A}} \mathbf{K}\) 在 K 上的秩）；若 \(0 \to M_{1} \to M_{2} \to M_{3} \to 0\) 是有限生成 A-模的正合列，则序列

\[
0 \rightarrow (\mathrm{M} _ {1}) _ {(\mathbf {K})} \rightarrow (\mathrm{M} _ {2}) _ {(\mathbf {K})} \rightarrow (\mathrm{M} _ {3}) _ {(\mathbf {K})} \rightarrow 0
\]

也正合，因此 \(r(\mathbf{M}_2) = r(\mathbf{M}_1) + r(\mathbf{M}_3)\) 。我们写

\[
\gamma (\mathbf {M}) = (r (\mathbf {M}), c (\mathbf {M})) \in \mathbf {Z} \times \mathbf {C} (\mathbf {A});
\]

于是 γ 满足命题 16 的性质 (i)，且当 M 是伪零模时 γ(M) = 0（因为 M 是扭模）。存在从 F(A) 到 Z × C(A) 的唯一映射，仍记作 γ，使得对每个有限生成 A-模 M 有 γ(M) = γ(cl(M))。我们将看到，上述性质本质上刻画了 γ：

命题 17. 设 G 是以加法写出的交换群，\(\phi\) 是从有限生成 A-模的类组成的集 \(F(A)\) 到 G 的映射；对每个有限生成 A-模 M，我们也滥用语言写 \(\phi(M) = \phi(cl(M))\) 。设下列条件满足：

\begin{enumerate}
\def\labelenumi{(\arabic{enumi})}
\item
  若 \(0 \to \mathbf{M}_1 \to \mathbf{M}_2 \to \mathbf{M}_3 \to 0\) 是有限生成 A-模的正合列，则 \(\phi(\mathbf{M}_2) = \phi(\mathbf{M}_1) + \phi(\mathbf{M}_3)\) 。
\item
  若 \(\mathbf{T}\) 是伪零的，则 \(\phi (\mathbf{T}) = 0\) 。
\end{enumerate}

则存在唯一同态 \(\theta: \mathbf{Z} \times \mathbf{C} \to \mathbf{G}\) 使 \(\phi = \theta \circ \gamma\) 。

由命题 16 (iv)，\(\mathbf{Z} \times \mathbf{C}\) 的每个元素对某个适当的有限生成 A-模 M 都形如 \((r(\mathbf{M}), c(\mathbf{M}))\) ；由此得 \(\theta\) 的唯一性。我们对 \(-\phi\) 在 \(T(\mathbf{A})\) 上的限制应用第 5 节的命题 11 ：于是存在同态 \(\theta_0: \mathbf{D} \to \mathbf{G}\) 使

\[
- \phi (\mathbf {T}) = \theta_ {0} (\chi (\mathbf {T}))
\]

对每个有限生成扭 A-模 T 成立。设 \(x\) 是 A 的非零元素；把性质 (1) 应用于正合列：

\[
0 \longrightarrow \mathrm{A} \xrightarrow {h _ {x}} \mathrm{A} \longrightarrow \mathrm{A} / x \mathrm{A} \longrightarrow 0
\]

其中 \(h_{x}\) 是乘以 x 的映射，我们得到 \(\phi(\mathbf{A}/x\mathbf{A}) = 0\) ，从而 \(\theta_{0}(\mathrm{div}(x)) = 0\) 。取商后，\(\theta_{0}\) 便定义了一个同态 \(\theta_{1}: C \to G\) ，且对每个扭 A-模 T 有 \(\phi(T) = \theta_{1}(c(T))\) 。我们现在证明同态 \(\theta\) （由 \(\theta(n, z) = n.\phi(\mathbf{A}) + \theta_{1}(z)\) 定义）解决了问题。为此，对每个有限生成 A-模 M 写 \(\phi'(\mathbf{M}) = \phi(\mathbf{M}) - \theta(\gamma(\mathbf{M}))\) ；显然当把 \(\phi\) 换成 \(\phi'\) 时条件 (1) 仍满足。此外，当 M 是扭模或自由模时 \(\phi'(\mathbf{M}) = 0\) ；但对每个有限生成 A-模 M，存在 M 的自由子模 L 使 M/L 是扭模（命题 15），故性质 (1) 表明对每个有限生成 A-模 M 都有 \(\phi'(\mathbf{M}) = 0\) 。

* 用 Abel 范畴的语言，命题 17 表明 \(\mathbf{Z} \times \mathrm{C}(\mathbf{A})\) 典范同构于商范畴 \(\mathcal{F}/\mathcal{F}'\) 的 Grothendieck 群，其中 \(\mathcal{F}\) 是有限生成 A-模的范畴，\(\mathcal{F}'\) 是 \(\mathcal{F}\) 中由伪零模组成的全子范畴。*

\subsection{关于纯量环有限扩张的性质}

在本节中，A 与 B 表示两个整闭 Noether 整环，满足 \(A \subset B\) ，B 是有限生成 A-模，而 K 与 L 分别是 A 与 B 的分式域。在涉及 A-模的场合，我们将用 \(div_{A}, \chi_{A}, c_{A}, \gamma_{A}, r_{A}\) 分别代替 div, \(\chi, c, \gamma, r\) ，对 B-模采用类似的记号。

我们知道（§ 1，第 10 节）：B 的素理想 P 高为 1，当且仅当 p = P ∩ A 高为 1；此外（出处同上，命题 14），对 p ∈ P(A)，位于 p 之上的素理想 P ∈ P(B) 只有有限多个。为简便起见，我们将用 P\textbar p 表示关系 ``P 位于 p 之上''（即 p = P ∩ A）；随后用 e\_\{P/p\} 或 e(P/p) 表示 v\_\{P\} 在 v\_\{p\} 之上的分歧指数 e(v\_\{P\}/v\_\{p\})（第 VI 章，§ 8，第 1 节），并用 f\_\{P/p\} 或 f(P/p) 表示剩余类次数 f(v\_\{P\}/v\_\{p\})（出处同上）；回忆离散赋值 v\_\{p\} 与 v\_\{P\} 都是赋范的，且 f\_\{P/p\} 是 B/P 的分式域在 A/p 的分式域上的次数。我们置 n = r\_A(B)，其中 B 被视为 A-模；于是按定义 n = {[}L: K{]}，且对一切 p ∈ P(A)，n 也是自由 A\_p-模 B 的秩，对一切 P\textbar p 皆然。于是由第 VI 章，§ 8，第 5 节，定理 2，对一切 p ∈ P(A)：

\[
\sum_ {\mathfrak {P} / \mathfrak {p}} e _ {\mathfrak {P} / \mathfrak {p}} f _ {\mathfrak {P} / \mathfrak {p}} = n.\tag{9}
\]

这样，由于 D(A) 与 D(B) 是自由 Z-模，我们用下列条件定义一个序群的递增同态 N: D(B) → D(A) （也记作 \(N_{B/A}\) ）：

\[
\mathrm{N} (\mathfrak {P}) = f _ {\mathfrak {P} / \mathfrak {p}}. \mathfrak {p} \quad \text {   for   } \quad \mathfrak {P} \in \mathrm{P} (\mathrm{B}), \quad \text {   where   } \quad \mathfrak {p} = \mathfrak {P} \cap \mathrm{A}.\tag{10}
\]

另一方面（§ 1，第 10 节，命题 14），我们已用下列条件定义了序群的递增同态 \(i\colon\mathbf{D}(\mathbf{A})\to\mathbf{D}(\mathbf{B})\) （也记作 \(i_{\mathbf{B}/\mathbf{A}}\) ）：

\begin{enumerate}
\def\labelenumi{(\arabic{enumi})}
\setcounter{enumi}{10}
\tightlist
\item
\end{enumerate}

\[
i (\mathfrak {p}) = \sum_ {\mathfrak {P} / \mathfrak {p}} e _ {\mathfrak {P} / \mathfrak {p}}. \mathfrak {P} \quad \text {   for   } \mathfrak {p} \in \mathrm{P} (\mathrm{A}).
\]

显然，对 A（相应地，B）的除子的每个族 \((d_{\iota})\) （相应地，\((d_{\iota}^{\prime}))\) ）：

\begin{enumerate}
\def\labelenumi{(\arabic{enumi})}
\setcounter{enumi}{11}
\tightlist
\item
\end{enumerate}

\[
\mathrm{N} (\sup (d _ {\iota} ^ {\prime})) = \sup (\mathrm{N} (d _ {\iota} ^ {\prime})), \quad \mathrm{N} (\inf (d _ {\iota} ^ {\prime})) = \inf (\mathrm{N} (d _ {\iota} ^ {\prime}))\tag{13}
\]

\[
i \left(\sup (d _ {\iota})\right) = \sup (i (d _ {\iota})), \quad i (\inf (d _ {\iota})) = \inf (i (d _ {\iota})).
\]

公式 (9) 表明：

\[
\mathbf {N} \circ i = n. 1 _ {\mathbf {D} (\mathbf {A})}.\tag{14}
\]

对一切 \(a \in \mathbf{A}\) （ \(\S 1\) ，第 10 节，命题 14）：

\[
i \left(\operatorname{div} _ {\mathbf {A}} (a)\right) = \operatorname{div} _ {\mathbf {B}} (a).\tag{15}
\]

我们（借助 (13)）推出，对 A 的每个分式理想 a，也有：

\[
i \left(\operatorname{div} _ {\mathbf {A}} (\mathfrak {a})\right) = \operatorname{div} _ {\mathbf {B}} (\mathfrak {a B}).\tag{16}
\]

对每个元素 \(b \in \mathbf{B}\) ，我们知道（第 V 章，§1，第 3 节，命题 11 的推论）\(\mathrm{N}_{\mathrm{L / K}}(b) \in \mathbf{A}\) ；此外（第 VI 章，§8，第 5 节，公式 (9)）：

\[
v _ {\mathfrak {p}} \left(\mathrm{N} _ {\mathrm{L} / \mathbf {K}} (b)\right) = \sum_ {\mathfrak {P} / \mathfrak {p}} f _ {\mathfrak {P} / \mathfrak {p}} v _ {\mathfrak {P}} (b)\tag{17}
\]

由此得：

\[
\mathbf {N} (\operatorname{div} _ {\mathbf {B}} (b)) = \operatorname{div} _ {\mathbf {A}} (\mathbf {N} _ {\mathbf {L} / \mathbf {K}} (b)).\tag{18}
\]

公式 (15) 与 (18) 表明，通过取商，同态 N 与 i 定义了一些同态，对它们我们将同样滥用记号，记作：

\[
\mathrm{N}: \mathrm{C} (\mathrm{B}) \rightarrow \mathrm{C} (\mathrm{A}), \quad i: \mathrm{C} (\mathrm{A}) \rightarrow \mathrm{C} (\mathrm{B}).
\]

注意，同态 \(i\colon\mathbf{C}(\mathbf{A})\to\mathbf{C}(\mathbf{B})\) 一般不是单射（§ 3，习题 7）。

回忆：对每个 B-模 R，\(R_{[A]}\) 表示把 R 的纯量限制到 A 而得到的 A-模（《代数》第 II 章，§1，第 13 节）。

命题 18. (i) R 是伪零 B-模，当且仅当 A-模 \(\mathbf{R}_{[\mathrm{A}]}\) 是伪零的。

\begin{enumerate}
\def\labelenumi{(\roman{enumi})}
\setcounter{enumi}{1}
\tightlist
\item
  R 是有限生成扭 B-模，当且仅当 \(R_{[A]}\) 是有限生成扭 A-模，且此时：
\end{enumerate}

\[
\chi_ {\mathbf {A}} (\mathbf {R} _ {[ \mathbf {A} ]}) = \mathbf {N} (\chi_ {\mathbf {B}} (\mathbf {R})).\tag{19}
\]

\begin{enumerate}
\def\labelenumi{(\roman{enumi})}
\setcounter{enumi}{2}
\tightlist
\item
  R 是有限生成 B-模，当且仅当 \(R_{[A]}\) 是有限生成 A-模，且此时：
\end{enumerate}

\begin{enumerate}
\def\labelenumi{(\arabic{enumi})}
\setcounter{enumi}{19}
\tightlist
\item
\end{enumerate}

\[
c _ {\mathbf {A}} (\mathrm{R} _ {[ \mathbf {A} ]}) = \mathrm{N} (c _ {\mathbf {B}} (\mathrm{R})) + r _ {\mathbf {B}} (\mathrm{R}) c _ {\mathbf {A}} (\mathrm{B})\tag{21}
\]

\[
r _ {\mathrm{A}} (\mathrm{R} _ {[ \mathrm{A} ]}) = n. r _ {\mathrm{B}} (\mathrm{R}) \quad (\text { recall   that } n = r _ {\mathrm{A}} (\mathrm{B})).
\]

由于 B 是有限生成 A-模，R 是有限生成 B-模当且仅当 \(R_{[A]}\) 是有限生成 A-模。此外，若 b 是 R 的零化子，则 \(b \cap A = a\) 是 \(R_{[A]}\) 的零化子；由于 B 在 A 上是整的，除 0 外没有理想位于 A 的理想 0 之上（第 V 章，§2，第 1 节，命题 1 的推论 1），故说 \(\mathfrak{a} \neq 0\) 与说 \(\mathfrak{b} \neq 0\) 是一回事。

\begin{enumerate}
\def\labelenumi{(\roman{enumi})}
\item
  依据这最后的注记，我们可以只考虑 R 是扭 B-模的情形。若 b 含于某个素理想 \(\mathfrak{P} \in \mathrm{P}(B)\) ，则 a 含于 \(\mathfrak{P} \cap A = p\) ，后者高为 1。反之，若 a 含于某个素理想 \(p \in \mathrm{P}(A)\) ，则存在 B 的素理想 P，它包含 b 且位于 p 之上（第 V 章，§2，第 1 节，定理 1 的推论 2）。断言 (i) 即由这些注记及第 4 节定义 2 得出。
\item
  对每个有限生成扭 B-模 R，我们写
\end{enumerate}

\[
\phi (\mathbf {R}) = \chi_ {\mathbf {A}} (\mathbf {R} _ {[ \mathbf {A} ]});
\]

显然（对有限生成扭 B-模）\(\phi\) 满足第 5 节命题 11 的条件 (1) 与 (2)（顾及 (i)）。因此存在同态 \(\theta: \mathrm{D}(\mathrm{B}) \to \mathrm{D}(\mathrm{A})\) ，使对每个有限生成扭 B-模 R 有 \(\phi(\mathrm{R}) = \theta(\chi_{\mathbf{B}}(\mathrm{R}))\) 。同态 \(\theta\) 由它在每个形如 \(\mathrm{B}/\mathfrak{P}\) （ \(\mathfrak{P} \in \mathrm{P}(\mathrm{B})\) ）的 B-模上的值确定，因为 \(\chi_{\mathbf{B}}(\mathrm{B}/\mathfrak{P}) = \mathfrak{P}\) 。现在，对每个素理想 \(\mathfrak{q} \neq \mathfrak{p} = \mathfrak{P} \cap \mathrm{A}\) （ \(\mathrm{P}(\mathrm{A})\) 中），有 \(\mathfrak{p} \notin \mathfrak{q}\) ，从而 \((\mathrm{B}/\mathfrak{P})_{\mathfrak{q}} = 0\) 。另一方面，若置 \(\mathrm{S} = \mathrm{A} - \mathfrak{p}\) ，则 \(\mathfrak{P} \cdot \mathrm{S}^{-1}\mathrm{B}\) 是 \(\mathrm{S}^{-1}\mathrm{B}\) 的极大理想，且 \((\mathrm{B}/\mathfrak{P})_{\mathfrak{p}} = \mathrm{S}^{-1}\mathrm{B}/\mathfrak{P}\) 。\(\mathrm{S}^{-1}\mathrm{B}\) 同构于 \(\mathrm{B}/\mathfrak{P}\) 的分式域（第 II 章，§ 2，第 5 节，命题 11），即赋值 \(v_{\mathfrak{P}}\) 的剩余域；因此它作为 \(\mathrm{A}_{\mathfrak{p}}\) -模的长度是 \(f_{\mathfrak{P}/\mathfrak{p}}\) ；这就证明了 \(\theta = \mathrm{N}\) （第 5 节，定义 4）。

\begin{enumerate}
\def\labelenumi{(\roman{enumi})}
\setcounter{enumi}{2}
\tightlist
\item
  若 T 是 R 的扭子模，则 \(T_{[A]}\) 是 \(R_{[A]}\) 的扭子模，且 \((\mathrm{R}/\mathrm{T})_{[\mathrm{A}]} = \mathrm{R}_{[\mathrm{A}]}/\mathrm{T}_{[\mathrm{A}]}\) ；因此为证 (21)，我们可以只考虑 R 无扭的情形。这时 R 等同于 \(\mathrm{R}_{(\mathrm{L})}\) 的一个子 B-模，且含有一组基 \((e_{i})_{1 \leqslant i \leqslant m}\) —— \(\mathrm{R}_{(\mathrm{L})}\) 在 L 上的基。若 \((b_{j})_{1 \leqslant j \leqslant n}\) 是 L 在 K 上的一组由 B 的元素组成的基，则诸 \(b_{j}e_{i}\) 构成 \(\mathrm{R}_{(\mathrm{L})}\) 在 K 上的一组由 R 的元素组成的基，由此得 (21)。另一方面，设 M 是 R 的自由子 B-模，使 R/M 是扭 B-模；由于 \(M_{[A]}\) 是 \(r_{\mathrm{B}}(\mathrm{R})\) 个同构于 B 的 A-模的直和，由命题 16 (i)
\end{enumerate}

\[
c _ {\mathbf {A}} (\mathbf {M} _ {[ \mathbf {A} ]}) = r _ {\mathbf {B}} (\mathbf {R}). c _ {\mathbf {A}} (\mathbf {B}).
\]

此外，依据 (19)，\(c_{\mathbf{A}}((\mathbf{R} / \mathbf{M})_{[\mathbf{A}]})=-c_{\mathbf{A}}(\mathbf{N}(\chi_{\mathbf{B}}(\mathbf{R} / \mathbf{M})))\) ；但由同态 \(\mathrm{N}\colon\mathrm{C}(\mathrm{B})\to\mathrm{C}(\mathrm{A})\) 的定义，\(c_{\mathbf{A}}(\mathrm{N}(d))=\mathrm{N}(c_{\mathbf{B}}(d))\) 对一切 \(d\in\mathrm{D}(\mathrm{B})\) 成立，且按定义有 \(c_{\mathbf{B}}(\chi(\mathrm{R}/\mathrm{M}))=-c_{\mathbf{B}}(\mathrm{R})\) ，最终得

\[
c _ {\mathbf {A}} ((\mathbf {R} / \mathbf {M}) _ {[ \mathbf {A} ]}) = \mathbf {N} (c _ {\mathbf {B}} (\mathbf {R}));
\]

于是只需应用命题 16 (i) 便得 (20)。

命题 19. 设 R 是有限生成 B-模。R 是自反的，当且仅当 \(R_{[A]}\) 是自反 A-模。

我们在命题 18 的证明中已指出，R 是无扭 B-模当且仅当 \(R_{[A]}\) 是无扭 A-模。因此可以设 R 是 \(W = R \otimes_{B} L\) 关于 B 的格。我们将使用下述引理：

引理 4. 设 W 是 L 上有限秩的向量空间，R 是 W 关于 B 的格。则对一切 \(\mathfrak{p} \in \mathrm{P}(\mathbf{A})\) ，\((\mathbf{R}_{[\mathbf{A}]})_{\mathfrak{p}} = \bigcap_{\mathfrak{P}/\mathfrak{p}} \mathbf{R}_{\mathfrak{P}}\) 。

若 \(S = A - p\) ，则环 \(S^{-1}B\) 的素理想由 \(B\) 中与 \(S\) 不相交的素理想生成，换言之即诸理想 \(\mathfrak{P}_i\) （ \(1 \leqslant i \leqslant m\) ，位于 \(p\) 之上）与理想 (0)；这表明 \(S^{-1}B\) 是半局部环，其极大理想是诸 \(m_i = \mathfrak{P}_i(S^{-1}B)\) （ \(1 \leqslant i \leqslant m\) ）；此外，局部环 \((S^{-1}B)_{\mathfrak{P}_i}\) 同构于 \(B_{m_i}\) （第 II 章，§ 2，第 5 节，命题 11），从而是离散赋值环。于是环 \(S^{-1}B\) 是 Dedekind 整环（§ 2，第 2 节，定理 1 (f)），且由于它是半局部的，它是主理想整环（§ 2，第 2 节，命题 1）。这样，\((R_{[A]})_p\) 就等于 \(S^{-1}R\) （视为 \(A_p\) -模） ；由上所述，\(S^{-1}R\) 是 \(W\) 关于 \(S^{-1}B\) 的自由格，故可对它应用第 2 节的定理 2，得 \(S^{-1}R = \bigcap_i (S^{-1}R)_{m_i}\) ；但 \((S^{-1}R)_{m_i} = R_{\mathfrak{P}_i}\) ，这就证明了引理。

回到命题 19 的证明，由引理 4，

\[
\bigcap_ {\mathfrak {P} \in P (B)} R _ {\mathfrak {P}} = \bigcap_ {\mathfrak {p} \in P (A)} \left(R _ {[ A ]}\right) _ {\mathfrak {p}}
\]

结论由第 2 节的定理 2 得出。

推论. 环 B 是自反 A-模。

命题 20. (i) 有限生成 A-模 M 是伪零的，当且仅当 \(\mathbf{M} \otimes_{\mathbf{A}} \mathbf{B}\) 是伪零 B-模。

\begin{enumerate}
\def\labelenumi{(\roman{enumi})}
\setcounter{enumi}{1}
\tightlist
\item
  若 \(\mathbf{M}\) 是有限生成扭 A-模，则 \(\mathbf{M} \otimes_{\mathbf{A}} \mathbf{B}\) 是有限生成 B-模，且：
\end{enumerate}

\[
\chi_ {\mathbf {B}} (\mathbf {M} \otimes_ {\mathbf {A}} \mathbf {B}) = i (\chi_ {\mathbf {A}} (\mathbf {M})).\tag{22}
\]

\begin{enumerate}
\def\labelenumi{(\roman{enumi})}
\setcounter{enumi}{2}
\tightlist
\item
  若 \(\mathbf{M}\) 是有限生成 A-模，则 \(\mathbf{M} \otimes_{\mathbf{A}} \mathbf{B}\) 是有限生成 B-模，且：
\end{enumerate}

\begin{enumerate}
\def\labelenumi{(\arabic{enumi})}
\setcounter{enumi}{22}
\tightlist
\item
\end{enumerate}

\[
c _ {\mathbf {B}} (\mathbf {M} \otimes_ {\mathbf {A}} \mathbf {B}) = i (c _ {\mathbf {A}} (\mathbf {M}))\tag{24}
\]

\[
r _ {\mathbf {B}} (\mathbf {M} \otimes_ {\mathbf {A}} \mathbf {B}) = r _ {\mathbf {A}} (\mathbf {M}).
\]

\begin{enumerate}
\def\labelenumi{(\roman{enumi})}
\tightlist
\item
  设 \(\mathfrak{P}\) 是 \(\mathbf{B},\mathfrak{p} = \mathfrak{P}\cap \mathbf{A}\) ；则 \((\mathbf{M}\otimes_{\mathbf{A}}\mathbf{B})_{\mathfrak{P}} = \mathbf{M}\otimes_{\mathbf{A}}\mathbf{B}_{\mathfrak{P}}\) （第 II 章，§ 2，第 7 节，命题 18），且另一方面
\end{enumerate}

\[
\mathbf {M} \otimes_ {\mathbf {A}} \mathbf {B} _ {\mathfrak {p}} = (\mathbf {M} \otimes_ {\mathbf {A}} \mathbf {A} _ {p}) \otimes_ {\mathbf {A} _ {p}} \mathbf {B} _ {\mathfrak {p}} = \mathbf {M} _ {p} \otimes_ {\mathbf {A} _ {p}} \mathbf {B} _ {\mathfrak {p}};
\]

因此关系 \(\mathbf{M}_{\mathfrak{p}} = 0\) 等价于 \((\mathbf{M} \otimes_{\mathbf{A}} \mathbf{B})_{\mathfrak{P}} = 0\) （第 II 章，§4，第 4 节，引理 4）。只需把这个注记应用于理想 \(\mathfrak{P} = (0)\) 以及诸理想 \(\mathfrak{P} \in \mathbf{P}(\mathbf{B})\) ，即可证明 (i)，其中用到第 4 节定义 2。

为证 (ii)，我们将使用下述引理：

引理 5. 设 \(\mathbf{M}_1, \mathbf{M}_2\) 是两个有限生成 A-模，\(f: \mathbf{M}_1 \to \mathbf{M}_2\) 是单同态。则 \(f \otimes 1_{\mathbf{B}}: \mathbf{M}_1 \otimes_{\mathbf{A}} \mathbf{B} \to \mathbf{M}_2 \otimes_{\mathbf{A}} \mathbf{B}\) 的核是伪零的。

设 \(\mathfrak{p}\) 是 A 的高 \(\leqslant 1\) 的素理想。则 \((\mathbf{M}_i\otimes_{\mathbf{A}}\mathbf{B})_{\mathfrak{p}} = (\mathbf{M}_i)_{\mathfrak{p}}\otimes_{\mathbf{A}_{\mathfrak{p}}}\mathbf{B}_{\mathfrak{p}}\) \((i = 1,2)\) （第 II 章，§2，第 7 节，命题 18），且 \((f\otimes 1_{\mathbf{B}})_{\mathfrak{p}} = f_{\mathfrak{p}}\otimes 1_{\mathbf{B}_{\mathfrak{p}}}\) ；\(f\) 是单射这一假设蕴含 \(f_{\mathfrak{p}}\) 也是单射（第 II 章，§2，第 4 节，定理 1）；另一方面，鉴于 \(\mathfrak{p},\mathbf{A}_{\mathfrak{p}}\) 是主理想整环，\(\mathbf{B}_{\mathfrak{p}}\) 是有限生成的无扭 \(\mathbf{A}_{\mathfrak{p}}\) -模，从而是自由的；我们断定 \(f_{\mathfrak{p}}\otimes 1_{\mathbf{B}_{\mathfrak{p}}}\) 本身是单射。若 \(\mathbf{I} = \operatorname {Ker}(f\otimes 1)\) ，则 \(\mathbf{I}_{\mathfrak{p}} = \operatorname {Ker}((f\otimes 1)_{\mathfrak{p}})\) （第 II 章，§2，第 4 节，定理 1）；因此 \(\mathrm{I_p} = 0\) ，故更有 \(\mathrm{I}_{\mathfrak{P}} = (\mathrm{I}_{\mathfrak{p}})_{\mathfrak{P}} = 0\) （对 \(\mathfrak{P}| \mathfrak{p}\) ），这就证明了引理（第 4 节，定义 2）。

现在回到 (ii) 的证明。对每个有限生成扭 A-模 M，写 \(\phi(M) = \chi_{B}(M \otimes_{A} B)\) ；由 (i) 可知，若 M 是伪零的，则 \(\phi(M) = 0\) 。另一方面，考虑有限生成扭 A-模的一个正合列：

\[
0 \rightarrow \mathrm{M} _ {1} \rightarrow \mathrm{M} _ {2} \rightarrow \mathrm{M} _ {3} \rightarrow 0.
\]

由引理 4 可知，存在 B-模的一个正合列：

\[
0 \rightarrow \mathrm{I} \rightarrow \mathrm{M} _ {1} \otimes_ {\mathrm{A}} \mathrm{B} \rightarrow \mathrm{M} _ {2} \otimes_ {\mathrm{A}} \mathrm{B} \rightarrow \mathrm{M} _ {3} \otimes_ {\mathrm{A}} \mathrm{B} \rightarrow 0
\]

其中 I 是伪零的。利用第 5 节命题 10 的推论，我们得到 \(\phi(\mathbf{M}_2) = \phi(\mathbf{M}_1) + \phi(\mathbf{M}_3)\) 。于是我们由第 5 节的命题 11 断定：存在同态 \(\theta: \mathbf{D}(\mathbf{A}) \to \mathbf{D}(\mathbf{B})\) ，使对每个有限生成扭 A-模 M 有 \(\phi(\mathbf{M}) = \theta(\chi_{\mathbf{A}}(\mathbf{M}))\) 。为证 \(\theta = i\) ，只需证明 \(\phi(\mathbf{A}/\mathfrak{p}) = i(\mathfrak{p})\) 对一切 \(\mathfrak{p} \in \mathbf{P}(\mathbf{A})\) 成立；现在 \((\mathbf{A}/\mathfrak{p}) \otimes_{\mathbf{A}} \mathbf{B} = \mathbf{B}/\mathfrak{p}\mathbf{B}\) ，且对一切 \(\mathfrak{P} \in \mathbf{P}(\mathbf{B})\) ，\((\mathbf{B}/\mathfrak{p}\mathbf{B})_{\mathfrak{P}} = \mathbf{B}_{\mathfrak{P}}/\mathfrak{p}\mathbf{B}_{\mathfrak{P}}\) ；若 \(\mathfrak{P}\) 不位于 \(\mathfrak{p}\) 之上，则最后的模是 0 ；反之，若 \(\mathfrak{P}| \mathfrak{p}, \mathbf{B}_{\mathfrak{P}}/\mathfrak{p}\mathbf{B}_{\mathfrak{P}}\) 是 \(\mathbf{B}_{\mathfrak{P}}\) -模，其长度为 \(e(\mathfrak{P}/\mathfrak{p})\) （按分歧指数的定义，第 VI 章，§ 8，第 1 节）；

因此 \(\chi_{\mathbf{B}}(\mathrm{B} / p\mathrm{B}) = \sum_{\mathfrak{P} / p}e_{\mathfrak{P} / p}\cdot \mathfrak{P} = i(p)\) ，这就完成了 (ii) 的证明。

公式 (24) 是直接的，因为

\[
(\mathbf {M} \otimes_ {\mathbf {A}} \mathbf {B}) \otimes_ {\mathbf {B}} \mathbf {L} = \mathbf {M} \otimes_ {\mathbf {A}} \mathbf {L} = (\mathbf {M} \otimes_ {\mathbf {A}} \mathbf {K}) \otimes_ {\mathbf {K}} \mathbf {L}
\]

且 \((\mathbf{M} \otimes_{\mathbf{A}} \mathbf{K}) \otimes_{\mathbf{K}} \mathbf{L}\) 在 L 上的秩等于 \(M \otimes_{A} K\) 在 K 上的秩。为证 (23)，考虑 M 的自由子模 H，使 Q = M/H 是扭 A-模。如上应用引理 4，我们得到 B-模的一个正合列：

\[
0 \rightarrow \mathrm{I} \rightarrow \mathrm{H} \otimes_ {\mathrm{A}} \mathrm{B} \rightarrow \mathrm{M} \otimes_ {\mathrm{A}} \mathrm{B} \rightarrow \mathrm{Q} \otimes_ {\mathrm{A}} \mathrm{B} \rightarrow 0
\]

因此由第 7 节命题 16 (ii)、(v) 以及命题 16 的推论 1 得

\[
c _ {\mathrm{B}} (\mathrm{M} \otimes_ {\mathrm{A}} \mathrm{B}) = c _ {\mathrm{B}} (\mathrm{Q} \otimes_ {\mathrm{A}} \mathrm{B}) = - c _ {\mathrm{B}} \left(\chi_ {\mathrm{B}} (\mathrm{Q} \otimes_ {\mathrm{A}} \mathrm{B})\right) = - c _ {\mathrm{B}} (i \left(\chi_ {\mathrm{A}} (\mathrm{Q})\right))
\]

这是依据 (ii)；但由同态 \(i: \mathbf{C}(\mathbf{A}) \to \mathbf{C}(\mathbf{B})\) 的定义，\(c_{\mathbf{B}}(i(\chi_{\mathbf{A}}(\mathbf{Q}))) = i(c_{\mathbf{A}}(\chi_{\mathbf{A}}(\mathbf{Q}))) = -i(c_{\mathbf{A}}(\mathbf{M}))\) ，这就完成了 (23) 的证明。

\textbf{注}\par

\begin{enumerate}
\def\labelenumi{(\arabic{enumi})}
\item
  若 M 是自反 A-模，则 \(M \otimes_{A} B\) 未必自反（习题 6）。但当 B 是平坦 A-模时它是自反的（第 2 节，命题 8）。
\item
  设 C 是第三个整闭 Noether 整环，满足 \(B \subset C\) ，且 C 是有限生成 B-模（从而也是有限生成 A-模）。则我们有传递公式：
\item
\end{enumerate}

\[
\mathbf {N} _ {\mathbf {C} / \mathbf {A}} = \mathbf {N} _ {\mathbf {B} / \mathbf {A}} \circ \mathbf {N} _ {\mathbf {C} / \mathbf {B}},\tag{26}
\]

\[
i _ {\mathbf {C} / \mathbf {A}} = i _ {\mathbf {C} / \mathbf {B}} \circ i _ {\mathbf {B} / \mathbf {A}}
\]

它们由分歧指数与剩余类次数的传递公式（第 VI 章，§ 8，第 1 节，引理 1）立即得出。

\subsection{约化定理}

我们保持第 2 节至第 7 节的记号与假设。

引理 6. 设 R 是交换环，\(\mathfrak{p}_i\) （ \(1 \leqslant i \leqslant n\) ）是 R 的互异素理想。

\begin{enumerate}
\def\labelenumi{(\roman{enumi})}
\item
  对 \(1 \leqslant i \leqslant n\) ，设 \(\mathbf{H}_i\) 是 \(\mathbf{R} / \mathfrak{p}_i\) 的满足下列条件的子集：不存在元素 \(\alpha_i \in \mathbf{R} / \mathfrak{p}_i\) ，使 \(\alpha_i + \mathbf{H}_i\) 含有某个 \(\neq 0\) 理想，即 \(\mathbf{R} / \mathfrak{p}_i\) 的一个理想。则存在 \(a \in \mathbf{R}\) ，使对 \(1 \leqslant i \leqslant n\) ，\(a\) 在 \(\mathbf{R} / \mathfrak{p}_i\) 中的典范像不属于 \(\mathbf{H}_i\) 。
\item
  若 \(\operatorname{Card}(\mathbf{H}_i) < \operatorname{Card}(\mathbf{R} / \mathfrak{p}_i)\) ，则诸 \(\mathbf{H}_i\) 满足 (i) 中的条件。
\item
  我们对 \(n\) 作归纳，\(n = 0\) 的情形是平凡的。于是设 \(n \geqslant 1\) 。我们可以对指标 \(i\) 作一个置换，从而可设 \(p_1\) 在诸 \(p_i\) 中极小，于是对 \(2 \leqslant i \leqslant n\) ，存在 \(c_i \in p_i\) 使 \(c_i \notin p_1\) 。由归纳假设，存在 \(b \in R\) ，使 \(b\) 在 \(R / p_i\) 中的典范像不属于 \(H_i\) （对 \(2 \leqslant i \leqslant n\) ）。对一切 \(x \in R\) ，我们写 \(a_{x}=b+xc_{2}c_{3}\ldots c_{n};\) 由于 \(c_{i}\in\mathfrak{p}_{i},\) 显然 \(a_{x}\equiv b\pmod{\mathfrak{p}_{i}}\) （对 \(2\leqslant i\leqslant n\) ）。因此只需证明存在 \(x\in R\) ，使 \(a_{x}\) 在 \(R/\mathfrak{p}_{1}\) 中的典范像不属于 \(H_{1}\) 。现在，当 x 跑遍 R 时，诸 \(a_{x}\) 在 \(R/\mathfrak{p}_{1}\) 中的典范像所成的集正是 \(\beta+c\) ，其中 \(\beta\) 是 b 的典范像，c 是 \(R/\mathfrak{p}_{1}\) 中由 \(c_{2}c_{3}\ldots c_{n}\) 的典范像生成的理想；鉴于诸 \(c_{i},c\neq0\) ，因为 \(R/\mathfrak{p}_{1}\) 是整环，而关于 \(H_{1}\) 的假设蕴含存在一个解决问题的 x。
\item
  由于 \(R/p_{i}\) 是整环，每个 \(\neq0\) 理想（\(R/p_{i}\) 中的）的基数都等于 \(R/p_{i}\) 的基数，且任一理想经 \(R/p_{i}\) 的元素所作的平移亦然，由此得结论。
\end{enumerate}

定理 6. 设 M 是无扭的有限生成 A-模。存在 M 的自由子模 L，使 M/L 同构于 A 的一个理想。

我们用 \(n\) 表示 \(M\) 的秩（即 \(V = M \otimes_{A} K\) 在 \(K\) 上的秩），并把 \(M\) 视为 \(V\) 关于 \(A\) 的格。于是对一切 \(p \in P(A)\) ，\(M_p\) 是 \(V\) 关于 \(A_p\) 的格（第 1 节，例 6），且由于 \(A_p\) 是主理想整环，\(M_p\) 是自由 \(A_p\) -模，秩为 \(n\) 。我们将写：

\[
\mathbf {M} (\mathfrak {p}) = \mathbf {M} _ {\mathfrak {p}} / \mathfrak {p} \mathbf {M} _ {\mathfrak {p}}.
\]

我们用 \(k(\mathfrak{p})\) 表示 \(\mathbf{A} / \mathfrak{p}\) 的分式域（同构于 \(\mathbf{A}_{\mathfrak{p}}\) 的剩余域）；于是 \(\mathbf{M}(\mathfrak{p}) = \mathbf{M} \otimes_{\mathbf{A}} k(\mathfrak{p})\) 是秩为 \(n\) 的向量空间（在 \(k(\mathfrak{p})\) 上）。对一切 \(x \in \mathbf{M}\) ，我们用 \(x(\mathfrak{p})\) 表示 \(x\) 在 \(\mathbf{M}(\mathfrak{p})\) 中的典范像。

引理 7. 设 \(x_{i}\) （ \(1 \leqslant i \leqslant m\) ）是 \(\mathbf{M}\) 中（在 \(\mathbf{A}\) 上或在 \(\mathbf{K}\) 上，二者等价）线性无关的元素，\(\mathbf{L}\) 是 \(\mathbf{M}\) 的由诸 \(x_{i}\) 生成的子 A-模。则对几乎所有 \(\mathfrak{p} \in \mathbf{P}\) ，\(x_{i}(\mathfrak{p}) \in \mathbf{M}(\mathfrak{p})\) 在 \(k(\mathfrak{p})\) 上线性无关；它们在 \(k(\mathfrak{p})\) 上对一切 \(\mathfrak{p} \in \mathbf{P}\) 线性无关，当且仅当 \(\mathbf{M}/\mathbf{L}\) 无扭。

设 \(x_{m+1}, \ldots, x_n\) 是 M 中与 \(x_1, \ldots, x_m\) 一起构成 V 的基的元素，N 是 M 的由诸 \(x_i (1 \leqslant i \leqslant n)\) 生成的自由子 A-模。由第 3 节的定理 3，\(N_p = M_p\) 对几乎所有 \(p \in P\) 成立；由于 \(x_1(p), \ldots, x_n(p)\) 构成 \(N(p)\) 在 \(k(p)\) 上的基，这就证明了第一个断言。

若 \(\mathbf{M} / \mathbf{L}\) 无扭，则 \((\mathbf{M} / \mathbf{L})_{\mathfrak{p}} = \mathbf{M}_{\mathfrak{p}} / \mathbf{L}_{\mathfrak{p}}\) 对一切 \(\mathfrak{p} \in \mathbb{P}\) 也无扭（第 1 节，例 6），且由于 \(\mathbf{A}_{\mathfrak{p}}\) 是主理想整环，\(\mathbf{M}_{\mathfrak{p}} / \mathbf{L}_{\mathfrak{p}}\) 是自由的。因此 \(\mathbf{M}_{\mathfrak{p}}\) 是 \(\mathbf{L}_{\mathfrak{p}}\) 与一个自由 \(\mathbf{A}_{\mathfrak{p}}\) -模 E （秩为 \(n - m\) ）的直和；于是 \(\mathbf{M}(\mathfrak{p})\) 是 \(\mathbf{L}(\mathfrak{p})\) 与 \(k(\mathfrak{p})\) -向量空间 \(\mathrm{E} / \mathfrak{p}\mathrm{E}\) （秩为 \(n - m\) ）的直和；因此 \(\mathbf{L}(\mathfrak{p})\) 的秩为 \(m\) ，而它由诸 \(x_{i}(\mathfrak{p})\) （ \(1 \leqslant i \leqslant m\) ）生成，故这些元素线性无关。

反之，设诸 \(x_{i}(\mathfrak{p})\) （ \(1 \leqslant i \leqslant m\) ）在 \(k(\mathfrak{p})\) 上对一切 \(\mathfrak{p} \in \mathbb{P}\) 线性无关。则 \(\mathbf{L}_{\mathfrak{p}}\) 是 \(\mathbf{M}_{\mathfrak{p}}\) 的直因子（对一切 \(\mathfrak{p}\) ）（第 II 章，§3，第 2 节，命题 5 的推论 1），因此 \(\mathbf{M}_{\mathfrak{p}} / \mathbf{L}_{\mathfrak{p}} = (\mathbf{M} / \mathbf{L})_{\mathfrak{p}}\) 对一切 \(p \in P\) 无扭。由第 IV 章，§1，第 2 节，命题 5 的推论，我们断定 \(\mathrm{P} \cap \mathrm{Ass}(\mathrm{M}/\mathrm{L}) = \varnothing\)。但由于 L 是自反的，由第 2 节命题 7 (i)，能属于 \(\mathrm{Ass}(\mathrm{M}/\mathrm{L})\) 的唯一素理想是理想 (0)；故 M/L 无扭。

引理 8. 设 M 的秩 n \(\geqslant2\) ；则存在 M 的元素 \(x \neq 0\) ，使 M/Ax 无扭。

设 \(y \neq 0\) 是 M 的一个元素。由引理 7，\(p \in P\) 中使 \(y(p) = 0\) 的元素所成的集 Y 是有限的。若 Y = \(\varnothing\) ，则由引理 7（应用于仅由单独一个元素 y 组成的序列 \((x_i)\) ）可知 M/Ay 无扭。因此设 Y ≠ \(\varnothing\) ，并写 S = \(\bigcap_{p \in Y} (A - p)\) ；我们知道（第 4 节，引理 2），\(S^{-1}A\) 是半局部主理想整环，其极大理想是诸 \(pS^{-1}A\) （p ∈ Y），相应的局部环是诸 \(A_{p}\) 。因此

\[
\mathrm{S} ^ {- 1} \mathrm{A} / \mathfrak {p} \mathrm{S} ^ {- 1} \mathrm{A} = k (\mathfrak {p}),
\]

由此得

\[
\begin{array}{c} \mathrm{S} ^ {- 1} \mathrm{M} / \mathfrak {p} \mathrm{S} ^ {- 1} \mathrm{M} = (\mathrm{M} / \mathfrak {p} \mathrm{M}) \otimes_ {\mathrm{A}} \mathrm{S} ^ {- 1} \mathrm{A} = \mathrm{M} \otimes_ {\mathrm{A}} ((\mathrm{A} / \mathfrak {p}) \otimes_ {\mathrm{A}} \mathrm{S} ^ {- 1} \mathrm{A}) = \\ \mathrm{M} \otimes_ {\mathrm{A}} k (\mathfrak {p}) = \mathrm{M} (\mathfrak {p}) \end{array}
\]

对一切 \(\mathfrak{p} \in \mathbf{Y}\) 。由第 II 章，§ 1，第 2 节，命题 6，存在元素 \(z / s \in \mathrm{S}^{-1}\mathbf{M}\) （ \(z \in \mathbf{M}, s \in \mathrm{S}\) ），它在诸 \(\mathbf{M}(\mathfrak{p})\) （ \(\mathfrak{p} \in \mathbf{Y}\) ）中的典范像都 \(\neq 0\) 。于是由 \(S\) 的定义，\(z(\mathfrak{p}) \neq 0\) 对一切 \(\mathfrak{p} \in Y\) 成立。我们还可以进一步设 \(y\) 与 \(z\) 在 \(K\) 上线性无关。因为在相反的情形，考虑元素 \(t \in M\) ，它与 \(y\) 线性无关（由于 \(n \geqslant 2\) ，这样的元素存在）；另一方面取一个元素 \(a \neq 0\) ，它属于 \(\bigcap_{\mathfrak{p} \in Y} \mathfrak{p}\) （由于 \(A\) 是整环，它不化为 0），并写 \(z' = z + at\) ：显然 \(y\) 与 \(z'\) 在 \(K\) 上线性无关，且 \(z'(\mathfrak{p}) = z(\mathfrak{p}) \neq 0\) 对一切 \(\mathfrak{p} \in Y\) 成立。

于是设 \(y\) 与 \(z\) 在 \(\mathbf{K}\) 上线性无关，令 \(\mathbf{Z}\) 为 \(\mathfrak{p} \in \mathbf{P} - \mathbf{Y}\) 中使 \(y(\mathfrak{p})\) 与 \(z(\mathfrak{p})\) 在 \(k(\mathfrak{p})\) 上线性无关者所成的集；由引理 7，这个集是有限的。对一切 \(\mathfrak{p} \in \mathbf{Z}\) ，我们可以写 \(z(\mathfrak{p}) = \lambda(\mathfrak{p})y(\mathfrak{p})\) ，其中 \(\lambda(\mathfrak{p}) \in k(\mathfrak{p})\) 。现在，\(\operatorname{Card}(\mathrm{A}/\mathfrak{p}) \geqslant 2\) 对一切 \(\mathfrak{p} \in \mathbf{P}\) 成立；因此由引理 6，存在 \(b \in \mathbf{A}\) ，使对一切 \(\mathfrak{p} \in \mathbf{Z}\) ，\(b\) 在 \(\mathrm{A}/\mathfrak{p}\) 中的典范像异于 \(\lambda(\mathfrak{p})\) 。我们现在证明元素 \(x = z - by\) 解决问题；只需（利用取 \(m = 1\) 的引理 7）验证 \(x(\mathfrak{p}) \neq 0\) 对一切 \(\mathfrak{p} \in \mathbf{P}\) 成立。现在：

--- 若 \(\mathfrak{p} \in \mathbf{Y}\) ，则按构造 \(x(\mathfrak{p}) \neq 0\) ；

--- 若 \(p \in Z\) ，则 \(x(p) = \mu_{p} y(p)\) ，其中由 b 的选取 \(\mu_{p} \neq 0\) ，故 \(x(p) \neq 0\) ，因为 \(y(p) \neq 0\) ；

- 若 \(\mathfrak{p} \in \mathrm{P} - (\mathrm{Y} \cup \mathrm{Z}), y(\mathfrak{p})\) 与 \(z(\mathfrak{p})\) 线性无关，故 \(x(\mathfrak{p}) \neq 0\) 。

这些引理既已建立，我们转向定理 6 的证明。我们对 n 作归纳，\(n \leqslant 1\) 的情形是平凡的，因为此时 M 本身同构于 A 的一个理想。设 \(n \geqslant 2\) ；由引理 8，存在 M 的秩 1 自由子模 \(L_{0}\) ，使 \(M/L_{0}\) 无扭；于是 \(M/L_{0}\) 的秩为 n - 1。由归纳假设，存在自由子模 \(L_{1}\) ，它含于 \(M/L_{0}\) ，使 \((\mathrm{M}/\mathrm{L}_{0})/\mathrm{L}_{1}\) 同构于 A 的一个理想。设 L 是 \(L_{1}\) 在 M 中的逆像；\(L/L_{0}\) 同构于 \(L_{1}\) ，且由于 \(L_{1}\) 自由，L 同构于 \(L_{0} \oplus L_{1}\) （《代数》第 II 章，§ 1，第 11 节，命题 21），从而是自由的；由于 M/L 同构于 \((\mathrm{M}/\mathrm{L}_{0})/\mathrm{L}_{1}\) ，定理得证。

注. 若 M 是自反的，M/L 未必自反（习题 9）。

\subsection{Dedekind 整环上的模}

我们现在假设 A 是 Dedekind 整环；于是我们知道理想 \(p \in P\) 都是极大的，且它们是 A 的仅有的 \(\neq 0\) 素理想（ \(\S 2\) ，第 1 节）；群 D(A) 等同于 A 的 \(\neq 0\) 分式理想所成的群 I(A)。

命题 21. 设 A 是 Dedekind 整环。每个伪零 A-模都是零。每个伪单射（相应地，伪满射、伪双射、伪零）A-模同态都是单射（相应地，满射、双射、零）。

第一个断言已证明（第 4 节，例 1）；其余的由它立即得出。

命题 22. 设 A 是 Dedekind 整环，M 是有限生成 A-模。下列性质等价：

\begin{enumerate}
\def\labelenumi{(\alph{enumi})}
\item
  M 无扭。
\item
  M 自反。
\item
  M 是投射的。
\end{enumerate}

我们已经知道（不对整环 A 作任何假设）：(b) 蕴含 (a)（第 2 节，注 1），(c) 蕴含 (b)（《代数》第 II 章，§ 2，第 7 节，命题 14 的推论 4）。若 M 无扭，它就等同于 \(\mathbf{V} = \mathbf{M} \otimes_{\mathbf{A}} \mathbf{K}\) 关于 A 的一个格；于是 \(\mathbf{M}_{\mathfrak{p}}\) 是自由 \(\mathbf{A}_{\mathfrak{p}}\) -模，对每个极大理想 \(\mathfrak{p} \in \mathbb{P}\) 皆然，因为 \(\mathbf{A}_{\mathfrak{p}}\) 是主理想整环。于是结论由第 II 章，§ 5，第 2 节，定理 1 (b) 得出。

推论. 设 M 是有限生成 A-模，T 是它的扭子模。则 T 是 M 的直因子。

由于 M/T 无扭且有限生成，由命题 22 它是投射的，于是推论由《代数》第 II 章，§ 2，第 2 节，命题 4 得出。

命题 23. 设 A 是 Dedekind 整环，T 是有限生成扭 A-模。存在两个有限族 \((n_{i})_{i\in I}\) 与 \((\mathfrak{p}_i)_{i\in I}\) ，其中诸 \(n_i\) 是 \(\geqslant 1\) 的整数，诸 \(\mathfrak{p}_i\) 是 P 的元素，使得 T 同构于直和 \(\bigoplus_{i\in I} (\mathrm{A} / \mathfrak{p}_{i}^{n_i})\) 。此外，族 \((n_{i})_{i\in I}\) 与 \((\mathfrak{p}_i)_{i\in I}\) 在相差指标集的一个双射的意义下是唯一的。

这由第 4 节的定理 5 得出，并注意到此处伪同构就是同构。

命题 24. 设 A 是 Dedekind 整环，M 是秩 \(n \geqslant 1\) 的有限生成无扭 A-模。则存在 A 的理想 \(d \neq 0\) ，使 M 同构于模 \(A^{n-1}\) 与 d 的直和。此外，理想 d 的类由这一条件唯一确定。

第 9 节的定理 6 表明存在 M 的自由子模 L，使 M/L 同构于 A 的理想 a。若 a = 0，我们取 d = A。否则，a 的秩为 1，于是 \(L = A^{n-1}\) 且 a 是投射模（命题 22）；因此 M 同构于 L 与 a 的直和（《代数》第 II 章，§2，第 2 节，命题 4），这就证明了命题的第一部分。此外，由第 7 节命题 16 (i)、(iv) 与 (v) 得 \(c(\mathbf{M}) = c(\mathfrak{d})\) ，由此得 d 的类的唯一性。

\textbf{注}\par

\begin{enumerate}
\def\labelenumi{(\arabic{enumi})}
\item
  命题 23、命题 24 以及命题 22 的推论完全确定了有限生成 A-模的结构。命题 24 表明，无扭的有限生成 A-模在同构意义下由它的秩及附加于它的除子类确定。
\item
  可以证明：在 Dedekind 整环上，非有限生成的投射模必然是自由的（习题 21），且投射模的每个子模都是投射的（习题 20）。
\end{enumerate}


\exercisesection{1}

\begin{enumerate}
\def\labelenumi{\arabic{enumi}.}
\tightlist
\item
  设 K 是域，A 是形式幂级数环 K{[}{[}T{]}{]} 中由 T 的系数为 0 的那些级数组成的子环；A 是一个 Noether 局部整环，但不是整闭的。
\end{enumerate}

\begin{enumerate}
\def\labelenumi{(\alph{enumi})}
\item
  证明 A 的每个非零且异于 A 的主理想都含有唯一的形如 \(T^{n} + \lambda T^{n+1}\) 的元素，其中 \(\lambda \in K\) ，\(n \geqslant 2\) ；此外，这个理想含有所有的幂 \(T^{m}\) ，其中 \(m \geqslant n + 2\) 。由此推出，A 的非主理想恰是诸理想 \(AT^{n} + AT^{n+1}\) ，其中 \(n \geqslant 2\) 。
\item
  证明 A 的每个分式理想都是除性的（证明每个非主理想是两个主理想之交）。描述幺半群 D(A) 的结构。A 的素理想只有极大理想 \(m = AT^{2} + AT^{3}\) 与 (0)。
\end{enumerate}

\begin{enumerate}
\def\labelenumi{\arabic{enumi}.}
\setcounter{enumi}{1}
\tightlist
\item
  \begin{enumerate}
  \def\labelenumii{(\alph{enumii})}
  \tightlist
  \item
    在域 K 上两个未定元的因子分解多项式整环 K{[}X, Y{]} 中，给出两个主理想 a, b 的例子，使 \(a + b\) 不是除性的。
  \end{enumerate}
\end{enumerate}

\begin{enumerate}
\def\labelenumi{(\alph{enumi})}
\setcounter{enumi}{1}
\item
  设 K 是域 \(\mathbf{Q}(\mathbf{X})\) （Q 上一个未定元的有理函数域）的二次扩张，它由对 \(\mathbf{Q}(\mathbf{X})\) 添加多项式 \(Y^{2}-2X\in\mathbf{Q}(\mathbf{X})[\mathbf{Y}]\) 的一个根 y 而得到。设 A 是 K 的由 Z、X 与 y 生成的子环；A 是整闭 Noether 整环。证明 \(p=AX+Ay\) 是高 1 的素理想，但 \(p^{2}\) 不是除性的（证明 \(X^{-1}p^{2}\) 是严格包含 p 的素理想）。此外 \(p.(A:p)\neq A\) 。
\item
  设 A 是整闭 Noether 局部整环，其极大理想 m 的高不是 \(\leqslant1\) ；设 p 是 A 的高 1 素理想；对每个整数 i \textgreater{} 0 ，令 \(a_{i} = p + m^{i}\) ；证明 \(\text{div}\left(\bigcap_{i} a_{i}\right)\) 异于 \(\text{sup}(\text{div}(a_{i}))\) 。
\end{enumerate}

\begin{enumerate}
\def\labelenumi{\arabic{enumi}.}
\setcounter{enumi}{2}
\tightlist
\item
  设 A 是 Noether 整环，p 是高 1 的素理想。
\end{enumerate}

  证明 A: \(p \neq A\) （在 A 的分式域中）。（若 \(c \neq 0\) 是 p 的元素，注意 \(p \in \text{Ass}(A/Ac)\) ，并据此推出 Ac: \(p \neq Ac\) ，依据是第 IV 章，§2，习题 30。）对非 Noether 的整环该结果仍成立吗（考虑一个高 1 的非离散赋值环）？

\begin{enumerate}
\def\labelenumi{\arabic{enumi}.}
\setcounter{enumi}{3}
\item
  设 A 是完全整闭整环，m 是 A 的一个极大理想。证明：或者 m 可逆，或者 m 不是除性的且 A: m = A。（注意，若 m 不可逆，则必有 m. \((\mathbf{A}:\mathfrak{m})^{k} = \mathfrak{m}\) 对每个整数 k \textgreater{} 0 成立。）
\item
  设 A 是整环，K 是它的分式域，U 是 \(K^{*}\) 的由 A 的可逆元组成的子群；把群 \(K^{*}/U = G\) 视为按如下关系有序的群，该关系由关系 \(x^{-1}y \in A\) （在 \(K^{*}\) 中）过渡到商而导出（《代数》第 VI 章，§ 1，第 5 节）。
\end{enumerate}

\begin{enumerate}
\def\labelenumi{(\alph{enumi})}
\item
  为使 \(\mathfrak{a} \neq 0\) （ \(\mathbf{K}\) 的一个分式理想）是除性的，当且仅当 \(\mathfrak{a} - \{0\}\) 在 \(\mathbf{G}\) 中的像是 \(\mathbf{G}\) 中的一个优势集（《代数》第 VI 章，§ 1，习题 30）；对每个分式理想 \(\mathfrak{a} \neq 0\) （ \(\mathbf{K}\) 的），\(\mathfrak{a} - \{0\}\) 的像是包含 \(\mathfrak{a} - \{0\}\) 的像的最小优势集。若 \(\mathbf{A}\) 是 Krull 整环，则 \(\mathbf{G}\) 是同构于形如 \(\mathbf{Z}^{(I)}\) 的有序群的一个定向子群的有序群。
\item
  证明：为使同态 \(w\) ——从 \(G\) 到全序群 \(\Gamma\) ——满足：它在 \(A - \{0\}\) 上对复合映射 \(K^* \to K^*/U \xrightarrow{w} \Gamma\) 的限制是一个赋值，当且仅当 \(w\) 是递增同态。
\item
  设 H 是乘积序群 \(Z \times Z\) 中由这样的有序对 \((s_{1}, s_{2})\) 组成的定向子群：它们满足 \(s_{1} + s_{2} \equiv 0 \pmod{2}\)。证明不存在整环 A，使 \(K^{*}/U\) 是同构于 H 的有序群（由 (b) 推出这样的整环必然是 Krull 整环，但对这个环第 5 节的命题 9 不成立（参看习题 22））。
\end{enumerate}

6. 设 A 是一个整环。\par

\begin{enumerate}
\def\labelenumi{(\alph{enumi})}
\item
  为使除性理想 \(\mathfrak{a}\) （ \(\mathbf{A}\) 中的）满足 \(\operatorname{div}(\mathfrak{a})\) 在 \(\mathbf{D}(\mathbf{A})\) 中可逆，当且仅当 \(\mathfrak{a} \colon \mathfrak{a} = \mathbf{A}\) （参看习题 5 及《代数》第 VI 章，§ 1，习题 30）。特别地，为使 \(\mathfrak{a} \colon \mathfrak{a} = \mathbf{A}\) 对每个除性理想 \(\mathfrak{a}\) （ \(\mathbf{A}\) 的）成立，当且仅当 \(\mathbf{A}\) 是完全整闭整环。
\item
  为使 \(\mathfrak{a} \colon \mathfrak{a} = \mathbf{A}\) 对每个有限生成理想 \(\mathfrak{a} \neq 0\) （ \(\mathbf{A}\) 的）成立，当且仅当 \(\mathbf{A}\) 是整闭的。
\end{enumerate}

\begin{enumerate}
\def\labelenumi{\arabic{enumi}.}
\setcounter{enumi}{6}
\item
  设 K 是域，\((v_{\iota})_{\iota \in I}\) 是 K 上满足第 3 节的条件 \((\mathrm{AK}_{\mathrm{I}})\) 与 \((\mathrm{AK}_{\mathrm{II}})\) 且满足下述条件的离散赋值族：对每个整数 r、每个由有理整数组成的族 \((n_{h})_{1 \leqslant h \leqslant r}\) 以及由 I 的互异元素组成的每个族 \((\iota_{h})_{1 \leqslant h \leqslant r}\) ，存在 \(x \in K\) ，使 \(v_{\iota_{h}}(x) = n_{h}\) 对 \(1 \leqslant h \leqslant r\) 成立，且 \(v_{\iota}(x) \geqslant 0\) 对 \(\iota\) （异于诸 \(\iota_{h}\) ）成立。证明诸 \(v_{\iota}\) 的赋值环之交是一个 Krull 整环 A，其分式域是 K，且 \((v_{\iota})_{\iota \in I}\) 是 A 的本性赋值族。（先证 K 是 A 的分式域，从而 A 是 Krull 整环；再证对一切 \(\iota \in I\) ，A 中由使 \(v_{\iota}(x) > 0\) 的 x 组成的素理想必然高为 1，用反证法论证并利用第 4 节命题 6 的推论 2。）
\item
  设 A 是 Krull 整环；证明对每个集合 I，多项式环 \(A[X_{i}]_{i\in I}\) 是 Krull 整环。（注意，若 B 是 Krull 整环，则 \(B[X_{1},\ldots,X_{n}]\) 的本性赋值在 B 上诱导的赋值就是 B 的本性赋值。）
\end{enumerate}

¶ 9. (a) 设 B 是离散赋值环，A = B{[}{[}X{]}{]} 是 B 上一个未定元的形式幂级数环，C = \(A_{x}\) 是分式环 \(S^{-1}A\)，其中 S 是由诸 \(X^{h}\) （h ≥ 0）组成的乘性子集。设 v 是 B 上的一个赋范赋值；对 C 的每个形如 f = ∑n≥h \(b_n X^{n}\) ≠ 0（\(b_h\) 在 B 中 ≠ 0，h 正或负）的元素，写 s(f) = v(\(b_h\) )。证明 s 是 C 上的欧几里得 stathm（《代数》第 VII 章，§ 1，习题 7），且对 C 中非零的 f, g 有 s(fg) = s(f) + s(g)。（若 s(f) = p ≥ s(g) = q 且 h 是 f 中 ≠0 项的诸次数中最小者，证明存在 u ∈ C 使 f - ug = \(X^{h+1} f_1\) ，其中 \(f_1\) ∈ C，并用反证法推出 C 中存在一种 ``欧几里得除法'' 过程。）若 s(f) = 0，则 f 在 C 中可逆。

* 证明 A 不是 Dedekind 整环。 *

\begin{enumerate}
\def\labelenumi{(\alph{enumi})}
\setcounter{enumi}{1}
\tightlist
\item
  由 (a) 推出：若 B 是 Krull 整环，则 A = B{[}{[}X{]}{]} 也是 Krull 整环。（若 K 是 B 的分式域，注意 A 是 K{[}{[}X{]}{]} 与一族 \((\mathrm{C}_{t})\) 与 A 有相同分式域的主理想整环之交，且 A 的每个元素在几乎所有 \(C_{t}\) 中可逆。）
\end{enumerate}

\begin{enumerate}
\def\labelenumi{\arabic{enumi}.}
\setcounter{enumi}{9}
\tightlist
\item
  设 A 是 Krull 整环，K 是它的分式域，L 是包含 K 的域，\((\mathbf{B}_{\alpha})\) 是包含 A 且含于 L 的 Krull 整环的一个右定向族，使得分式域 \(L_{\alpha}\) （ \(B_{\alpha}\) 的）是 K 的有限代数扩张，且 \(B_{\alpha}\) 是 A 在 \(L_{\alpha}\) 中的整闭包。对 A 的每个本性赋值 v 与每个 \(\alpha\) ，令 \(e_{\alpha}(v)\) 为 \(B_{\alpha}\) 上扩张 v 的诸赋值的分歧指数之和。证明：为使诸 \(B_{\alpha}\) 之并是 Krull 整环，当且仅当对 A 的每个本性赋值 v，诸 \(e_{\alpha}(v)\) 组成的集有界。
\end{enumerate}

¶ 11. 除子 \(d\) （在整环 \(\mathbf{A}\) 中）称为有限生成的，如果它形如 \(\operatorname{div}(a)\) ，其中 \(a\) 是一个有限生成分式理想（它未必是除性的；参看 (b)）。

\begin{enumerate}
\def\labelenumi{(\alph{enumi})}
\item
    证明：若 A 是 Krull 整环，则 D(A) 的每个除子都形如 div(a)，其中 \(a = Ax + Ay\) ，x、y 是 A 的分式域中 \(\neq 0\) 的元素（利用第 5 节命题 9 的推论 2）。
\item
  设 K 是域，\(A = K[X_{n}]_{n \in N}\) 是 K 上可数多个未定元的多项式环，它是 Krull 环（习题 8）。设 \(A'\) 是 A 的由单位元与单项式 \(X_{i}X_{j}\) （对所有对 \((i,j)\) ）生成的子环（等价地说，\(A'\) 是所有各项次数均为偶数的多项式组成的集）。若 L 与 \(L'\) 是 A 与 \(A'\) 的分式域，证明 \(A' = A \cap L'\) ，从而 \(A'\) 是 Krull 环。设 p 是 \(A'\) 中由乘积 \(X_{0}X_{i}\) （对所有 \(i \geqslant 0\) ）生成的理想；证明 p 是高 1 的素理想（从而是除性的），但它不是有限生成的。
\end{enumerate}

\begin{enumerate}
\def\labelenumi{\arabic{enumi}.}
\setcounter{enumi}{11}
\tightlist
\item
  \begin{enumerate}
  \def\labelenumii{(\alph{enumii})}
  \tightlist
  \item
    设 \(\mathbf{A}\) 是整闭整环，\(f(\mathbf{X}) = \sum_{i} a_i \mathbf{X}^i\) 与 \(g(\mathbf{X}) = \sum_{j} b_j \mathbf{X}^j\) 是 \(\mathbf{A}[\mathbf{X}]\) 中的两个多项式；证明：若 \(c \in \mathbf{A}\) 使 \(f(\mathbf{X}) g(\mathbf{X})\) 的所有系数都属于 \(\mathbf{A}c\) ，则所有乘积 \(a_i b_j\) 都属于 \(\mathbf{A}c\) （化归为 \(\mathbf{A}\) 是赋值环的情形）。
  \end{enumerate}
\end{enumerate}

\begin{enumerate}
\def\labelenumi{(\alph{enumi})}
\setcounter{enumi}{1}
\item
  若 A 是习题 1 中定义的环，给出两个次数为 1 的多项式 \(f, g\) （在 A{[}X{]} 中）及一个元素 \(c \in \mathbf{A}\) 的例子，使 (a) 的结论对它们不成立（取 \(c = \mathbf{T}^4 + \mathbf{T}^5\) ）。
\item
  在 (a) 的假设下，设 a、b、c 分别是由 f、g 与 fg 的系数生成的 A 的理想：证明 \(\text{div}(c) = \text{div}(a) + \text{div}(b)\) 。给出 \(\tilde{c} \neq \tilde{a}\tilde{b}\) 的一个例子。证明：若 c 是主理想，则 c = ab，且 a 与 b 都可逆；若进而 A 是局部环，则 a 与 b 都是主理想。
\end{enumerate}

\begin{enumerate}
\def\labelenumi{\arabic{enumi}.}
\setcounter{enumi}{12}
\tightlist
\item
  设 A 是整环，\((p_{i})_{i\in I}\) 是一族非零元素，使得诸主理想 \(Ap_{i}\) 都是素理想；设 S 是由诸 \(p_{i}\) 生成的乘性子集。
\end{enumerate}

\begin{enumerate}
\def\labelenumi{(\alph{enumi})}
\item
  证明 \(\mathbf{A} = \mathbf{S}^{-1}\mathbf{A}\cap \left(\bigcap_{i}\mathbf{A}_{\mathbf{A}p_i}\right)\) 。
\item
  假设 A 的主理想的每个非空族都有极大元。证明每个环 \(A_{Ap_{t}}\) 都是离散赋值环（参看第 VI 章，§3，第 5 节，命题 9）。由此推出：若 \(S^{-1}A\) 是整闭的（相应地，完全整闭的），则 A 是整闭的（相应地，完全整闭的）。若 \(S^{-1}A\) 是 Krull 整环，证明 A 是 Krull 整环。
\end{enumerate}

\begin{enumerate}
\def\labelenumi{\arabic{enumi}.}
\setcounter{enumi}{13}
\tightlist
\item
  设 A 是 Krull 整环，S 是 A 的一个不含 0 的饱和乘性子集（《代数》第 II 章，§ 2，习题 1）。证明：若典范同态 \(\bar{i}\colon\mathrm{C}(\mathrm{A})\to\mathrm{C}(\mathrm{S}^{-1}\mathrm{A})\) （第 10 节，命题 17）是双射，则 S 是由
\end{enumerate}
